\documentclass[11pt]{article}
\usepackage[margin=0.95in]{geometry}
\usepackage{amsmath,amssymb,amsthm,booktabs}
\usepackage[colorlinks=true,linkcolor=blue,citecolor=blue,urlcolor=blue]{hyperref}
\newtheorem{theorem}{Theorem}[section]
\newtheorem{lemma}[theorem]{Lemma}
\newtheorem{corollary}[theorem]{Corollary}
\theoremstyle{remark}\newtheorem{remark}[theorem]{Remark}
\newcommand{\F}{\mathbb F}
\newcommand{\Z}{\mathbb Z}
\newcommand{\Mat}{\operatorname{Mat}}
\newcommand{\GL}{\operatorname{GL}}
\newcommand{\tr}{\operatorname{tr}}
\newcommand{\im}{\operatorname{im}}
\newcommand{\rk}{\operatorname{rank}}
\newcommand{\ir}{\operatorname{ir}}
\newcommand{\Id}{\operatorname{Id}}
\title{Sharp arity bounds for clonoids from mixed-exponent groups}
\author{Henry Zweiman}
\date{September 9, 2026}
\begin{document}
\maketitle
\begin{abstract}
Let $p$ be a prime and $s\ge0$. We prove that every clonoid from
$\Z/p^2\Z\oplus(\Z/p\Z)^s$ to a module on which $p$ is invertible is generated by its $(s+2)$-ary functions. The stronger uniform reconstruction identity holds over $\Z[1/p]$, and the bound is optimal in general, even for finite field targets. In particular, this answers the mixed-exponent case singled out by Fioravanti, Kompatscher and Rossi. The proof combines Kov\'acs's theorem on finite-field matrix semigroup algebras with the character decomposition of a principal congruence subgroup. A correction on the $p$-torsion submodule handles the zero character. A Fourier functional gives the matching lower bound for actual clonoid generation.
\end{abstract}

\section{Introduction and main result}

For algebras $A$ and $B$, a clonoid from $A$ to $B$ is a collection of functions $A^n\to B$, for positive integers $n$, closed under precomposition by term operations of $A$ and postcomposition by term operations of $B$. When both algebras are modules, these operations are linear combinations over their respective scalar rings. We regard an abelian group as a module over the integers.

Mayr and Wynne \cite{MW} proved finiteness of the clonoid lattice for coprime finite modules when the source has a distributive submodule lattice. Fioravanti, Kompatscher and Rossi \cite{FKR} established sharp uniform arity bounds for finite vector-space sources. Their Conjecture 1 predicts finiteness for all pairs of finite modules of coprime order. Their stronger Conjecture 59 asks for uniform reconstruction for finite abelian $p$-groups, and Question 60 asks specifically about $\Z/p^2\Z\oplus\Z/p\Z$.

We settle that question and the family obtained by adjoining any number of elementary abelian factors. The case $s=0$ below recovers a known cyclic case of \cite{MW}; the mixed-exponent cases have $s\ge1$.

For a matrix $M$ over a ring $R$, its \emph{inner rank} $\ir_R(M)$ is the least $k\ge0$ for which $M=UV$ with $U$ having $k$ columns and $V$ having $k$ rows. The zero matrix has inner rank zero. A substitution $x\mapsto Mx$ with inner rank at most $k$ factors through $k$ variables.

\begin{theorem}\label{thm:main}
Let $p$ be a prime, $s\ge0$, and
\[
 A=\Z/p^2\Z\oplus(\Z/p\Z)^s,\qquad D=\Z[1/p].
\]
For every $n\ge1$ there are coefficients $a_M\in D$, indexed by matrices $M\in\Mat_n(\Z/p^2\Z)$ of inner rank at most $s+2$, such that, for every $D$-module $B$ and every function $f:A^n\to B$,
\begin{equation}\label{eq:uniform}
 f(x)=\sum_{\ir(M)\le s+2}a_M f(Mx)\qquad(x\in A^n).
\end{equation}
Consequently every clonoid from $A$ to a module whose scalar ring contains $p^{-1}$ is generated by its $(s+2)$-ary part.

For every $p,s$ there is a finite field $k$ of characteristic different from $p$ and a clonoid from $A$ to the regular $k$-module that is not generated by its $(s+1)$-ary part.
\end{theorem}

The matrix in \eqref{eq:uniform} acts by integer scalar combinations on the whole group $A$. Its actions on the cyclic and elementary factors are compatible reductions of the same matrix. The proof does not assume that one can choose these actions independently.

\begin{corollary}\label{cor:finite}
If $B$ is a finite module, there are finitely many clonoids from $A$ to $B$ if and only if $p$ does not divide $|B|$.
\end{corollary}
\begin{proof}
If $p\nmid|B|$, multiplication by $p$ has an inverse given by an integer scalar on $B$. Thus \eqref{eq:uniform} uses term operations of $B$. Each clonoid is determined by its $(s+2)$-ary part, a subset of the finite set $B^{A^{s+2}}$. Conversely, when $p\mid|B|$, the common-prime-divisor result of Mayr and Wynne \cite[Theorem 1.3]{MW} gives an infinite ascending chain of clonoids.
\end{proof}

\section{Matrix operators and the finite-field input}

If $X$ is a finite set, write $K[X]$ for the free module with basis $[x]$, $x\in X$, over a commutative ring $K$. A map $g:X\to X$ induces the pushforward $[x]\mapsto[g(x)]$. Transposing an identity between these operators gives the corresponding identity for precomposition on functions.

We use the following consequence of Kov\'acs's theorem \cite{Ko}. It is also the operator form of the finite-vector-space reconstruction theorem in \cite[Theorem 49]{FKR}. We include a derivation from the older matrix-algebra result.

\begin{lemma}\label{lem:field}
Let $F$ be a finite field of characteristic $p$, and let $K$ be a commutative ring in which $p$ is invertible. For integers $m,d\ge0$, the identity on $K[F^{m\times d}]$ is a $K$-linear combination of operators
\[
 [X]\longmapsto[CX],\qquad C\in\Mat_m(F),\quad \rk(C)\le d.
\]
\end{lemma}
\begin{proof}
If $d=0$, the set has one element and the zero matrix suffices. If $m\le d$, use the identity matrix. For $m>d\ge1$, Kov\'acs's theorem says that the span of singular matrices in $K[\Mat_m(F)]$ has its own multiplicative identity $e$; see also \cite[Theorem 2.1]{Kuhn}. For each $X\in F^{m\times d}$ choose a singular projection $P$ whose image contains the column space of $X$. Since $e[P]=[P]$, its action gives
\[
 e[X]=e[P][X]=[P][X]=[X].
\]
Thus $e$ acts as the identity. Expand $e$ in singular matrices. Factor each such matrix as $UW$ through $F^r$, where $r<m$. By induction, insert a rank-at-most-$d$ operator identity on $K[F^{r\times d}]$ between the maps induced by $W$ and $U$. Each resulting matrix has the form $UCW$ and rank at most $d$. This proves the assertion by induction on $m$.
\end{proof}

We now fix $n$ and put
\[
 R=\Z/p^2\Z,\quad F=\F_p,\quad
 K=D[\zeta],\qquad \zeta=\zeta_p.
\]
Here $K$ is the cyclotomic extension of $D$ with basis
$1,\zeta,\ldots,\zeta^{p-2}$. Let
\[
 V=K[A^n]=K[R^n\times F^{n\times s}].
\]
Write its basis elements as $[u,v]$, where the columns of $v$ record the elementary coordinates. For $M\in\Mat_n(R)$ define
\begin{equation}\label{eq:rho}
 \rho(M)[u,v]=[Mu,\overline Mv].
\end{equation}
A bar denotes reduction modulo $p$. Let $\mathcal E$ be the image of $K[\Mat_n(R)]$ in $\operatorname{End}_K(V)$, and set
\[
 \mathcal I_j=\operatorname{span}_K\{\rho(M):\ir_R(M)\le j\}.
\]
This is a two-sided ideal of $\mathcal E$, since multiplication of matrices cannot increase inner rank. Our task is to show $\Id_V\in\mathcal I_{s+2}$.

Whenever a field matrix $C$ of rank at most $d$ is lifted to $R$, we choose a factorization $C=UW$ through $F^d$, lift $U,W$ entrywise, and use $\widetilde U\widetilde W$. This produces a lift of inner rank at most $d$. An arbitrary entrywise lift of $C$ need not have that property.

\section{The nonzero congruence characters}

The principal congruence subgroup
\[
 \Gamma=\{I+pH:H\in\Mat_n(F)\}\subseteq\GL_n(R)
\]
is elementary abelian. The expression $pH$ is independent of the chosen lifts of the entries of $H$. For $T\in\Mat_n(F)$ put
\begin{equation}\label{eq:projector}
 e_T=p^{-n^2}\sum_{H\in\Mat_n(F)}
       \zeta^{-\tr(TH)}\rho(I+pH).
\end{equation}
Orthogonality of the characters of the additive matrix group gives
\begin{equation}\label{eq:orthogonal}
 e_Te_{T'}=0\ (T\ne T'),\qquad e_T^2=e_T,
 \qquad\sum_Te_T=\Id_V.
\end{equation}
These identities hold over $K$: all normalizing denominators are powers of $p$, and the sum of a nontrivial additive character is zero.

\begin{lemma}\label{lem:rankone}
If $\rk(T)>1$, then $e_T=0$ on $V$. If $T\ne0$ has rank one, then $e_T\in\mathcal I_{s+2}$.
\end{lemma}
\begin{proof}
Fix a residue vector $a=\overline u$ and an elementary matrix $v$. If $a=0$, then $\Gamma$ fixes every $[u,v]$ with this residue. If $a\ne0$, its orbit consists of
\[
 [u+pw,v]\qquad(w\in F^n),
\]
and the stabilizer is $Ha=0$. The annihilator of this stabilizer under the trace pairing consists precisely of matrices $T=ab^t$, $b\in F^n$: indeed $\tr(ab^tH)=b^tHa$, and both spaces have dimension $n$. Therefore only the zero character and rank-one characters occur.

We can assume $n>s+2$, since otherwise every $e_T$ already belongs to $\mathcal I_{s+2}=\mathcal E$. A nonzero rank-one matrix is conjugate to one of
\[
 T=tE_{11}\quad(t\in F^*),\qquad T=E_{12}.
\]
For the first case use its nonzero image eigenvector and its kernel; for the second use a vector in the image and a preimage of that vector. A conjugating matrix over $F$ lifts to $\GL_n(R)$. Conjugating \eqref{eq:projector} correspondingly conjugates $T$ and preserves $\mathcal I_{s+2}$, so it is enough to handle these forms.

Set $h=1$ and $b^t=te_1^t$ in the first case, and $h=2$ and $b^t=e_2^t$ in the second. For $c\in F^*$ and $v\in F^{n\times s}$, use the integer representative of $c$ in $R$ and define
\begin{equation}\label{eq:Ecv}
 E_{c,v}=p^{-n}\sum_{w\in F^n}
    \zeta^{-c^{-1}b^tw}[ce_1+pw,v].
\end{equation}
These vectors form a basis of $e_TV$. To see this, the only possible nonzero residues are $ce_1$, and each such orbit contributes the one-dimensional character summand displayed in \eqref{eq:Ecv}. In particular,
\begin{equation}\label{eq:phase}
 e_T[ce_1+pz,v]=\zeta^{c^{-1}b^tz}E_{c,v}.
\end{equation}

Take $C\in\Mat_{n-h}(F)$ and a factor-preserving lift $\widetilde C$ as above. Put
\[
 M_C=\operatorname{diag}(I_h,\widetilde C).
\]
It fixes $ce_1$ exactly and satisfies $b^t\overline M_C=b^t$. Expanding \eqref{eq:Ecv} and then using \eqref{eq:phase} gives
\begin{equation}\label{eq:sandwich}
 e_T\rho(M_C)e_TE_{c,v}=E_{c,\overline M_Cv}.
\end{equation}
Indeed, the summand indexed by $w$ acquires the phase
$\zeta^{c^{-1}b^t\overline M_Cw}$, which cancels its coefficient phase. All $p^n$ summands then agree. This calculation does not assume that $C$ is invertible.

Apply Lemma \ref{lem:field} to the last $n-h$ rows of $v$, which have $s$ columns, while keeping its first $h$ rows and $c$ fixed. If $\alpha_C$ are the resulting coefficients with $\rk(C)\le s$, then \eqref{eq:sandwich} yields the operator identity
\begin{equation}\label{eq:nonzero}
 e_T=\sum_C\alpha_Ce_T\rho(M_C)e_T.
\end{equation}
Each $M_C$ has inner rank at most $h+s\le s+2$. Since $\mathcal I_{s+2}$ is a two-sided ideal, the right hand side belongs to it. This proves the lemma, including $s=0$ by the singleton case of Lemma \ref{lem:field}.
\end{proof}

\section{The zero character and coefficient descent}

Let $W\subset V$ be the span of the basis elements with $\overline u=0$. This is an $\mathcal E$-submodule. By writing $u=pw$, we identify it with
\begin{equation}\label{eq:torsion}
 W=K[F^{n\times(s+1)}],
\end{equation}
on which $\rho(M)$ acts through $\overline M$.

The image $e_0V$ contains $W$. Modulo $W$ it has basis
\begin{equation}\label{eq:zero-basis}
 E_{a,v}=p^{-n}\sum_{w\in F^n}[\widetilde a+pw,v]+W,
 \qquad a\in F^n\setminus\{0\},
\end{equation}
where the choice of the lift $\widetilde a$ is immaterial. For any $M\in\Mat_n(R)$, averaging the output orbit shows
\begin{equation}\label{eq:zero-action}
 e_0\rho(M)e_0E_{a,v}=
 \begin{cases}
 E_{\overline Ma,\overline Mv},&\overline Ma\ne0,\\
 0,&\overline Ma=0,
 \end{cases}
 \quad\text{in }V/W.
\end{equation}
For the first case, $e_0$ sends every output basis element to the average of all lifts of the same nonzero residue, so multiplicities sum to one. This also removes the carry in $M\widetilde a$. For the second case, all output basis elements lie in $W$.

Thus the compressed action in \eqref{eq:zero-action} is the field action on
$K[F^{n\times(s+1)}]$ modulo the invariant subspace in which the first column is zero. Choose a field identity from Lemma \ref{lem:field}, with coefficients $\beta_C$ and $\rk(C)\le s+1$, and factor-preserving lifts $M_C$. Define
\begin{equation}\label{eq:B-L-J}
 B_0=\sum_C\beta_C\rho(M_C),\qquad
 L=e_0B_0e_0,\qquad J=L+\sum_{T\ne0}e_T.
\end{equation}
Here $B_0,L\in\mathcal I_{s+1}$, and $J\in\mathcal I_{s+2}$ by Lemma \ref{lem:rankone}. Equation \eqref{eq:torsion} says that $B_0$ is the identity on $W$. Equations \eqref{eq:orthogonal} and \eqref{eq:zero-action} say that $(\Id_V-J)V\subseteq W$. Consequently
\begin{equation}\label{eq:correction}
 \Id_V=J+B_0(\Id_V-J)\in\mathcal I_{s+2}.
\end{equation}
This proves the desired operator identity over $K$.

All matrices of the operators $\rho(M)$ in the basis $[u,v]$ have integer entries. Apply to their coefficients the $D$-linear map
\[
 K\longrightarrow D,\qquad
 \sum_{i=0}^{p-2}d_i\zeta^i\longmapsto d_0.
\]
It fixes $1$ and therefore descends the identity to $D$. No averaging over a group of order $p-1$ is involved. Transposing now gives \eqref{eq:uniform} for $D$-valued functions, and the same finite coefficient identity holds for functions into every $D$-module.

To obtain clonoid generation, write each matrix in \eqref{eq:uniform} as $M=UV$ through $s+2$ variables. If $f$ belongs to a clonoid $\mathcal C$, then $f\circ U$ belongs to its $(s+2)$-ary part, and $f\circ U\circ V$ is generated by that part. The scalar coefficients are term operations on the target module. This proves the upper bound in Theorem \ref{thm:main}.

\begin{remark}
The zero-character image need not be stable under singular substitutions. The proof uses the sandwiches $e_0\rho(M)e_0$ and the invariant submodule $W$ explicitly. Replacing \eqref{eq:zero-action} by an unqualified identification with the whole residue-field representation would overlook its torsion contribution.
\end{remark}

\section{A sharp lower bound}

Choose a finite field $k$ of characteristic different from $p$ containing a primitive $p$th root $\zeta$. For example, for any prime $\ell\ne p$ one may use $k=\F_{\ell^{p-1}}$. Set $n=s+2$ and
\[
 V_0=[e_3\ \cdots\ e_n]\in F^{n\times s};
\]
for $s=0$ this matrix has no columns. Define $f:A^n\to k$ by
\begin{equation}\label{eq:lower-function}
 f(e_1+pw,V_0)=\zeta^{w_2}\quad(w\in F^n),
 \qquad f=0\text{ elsewhere}.
\end{equation}
On the space of all $k$-valued functions on $A^n$ define
\begin{equation}\label{eq:functional}
 \Lambda(g)=p^{-n}\sum_{w\in F^n}
                \zeta^{-w_2}g(e_1+pw,V_0).
\end{equation}
Then $\Lambda(f)=1$.

\begin{lemma}\label{lem:lower}
If $M\in\Mat_n(R)$ has inner rank at most $n-1$, then
$\Lambda(f\circ M)=0$.
\end{lemma}
\begin{proof}
For any summand in \eqref{eq:functional} to be nonzero, we must have
\begin{equation}\label{eq:fixed-columns}
 \overline Me_1=e_1,\qquad
 \overline Me_j=e_j\quad(3\le j\le n).
\end{equation}
If these conditions fail, every summand vanishes. If they hold, write $Me_1=e_1+pz$. Substitution in \eqref{eq:lower-function} gives
\[
 \Lambda(f\circ M)=\zeta^{z_2}p^{-n}
       \sum_{w\in F^n}\zeta^{(e_2^t\overline M-e_2^t)w}.
\]
This is zero unless $e_2^t\overline M=e_2^t$. Together with \eqref{eq:fixed-columns}, that equality forces $\overline M$ to be invertible: all columns other than the second are the corresponding standard basis vectors, and the second coordinate of the second column is one. But $\rk(\overline M)\le\ir_R(M)\le n-1$, a contradiction.
\end{proof}

Let $\mathcal C$ be the clonoid generated by $f$. Every member of its $r$-ary part is a $k$-linear combination of substitutions of $f$ by matrices with $r$ input columns. Hence every $n$-ary function generated by $\mathcal C^{(n-1)}$ is a $k$-linear combination of functions $f\circ M$ with $\ir_R(M)\le n-1$. This follows directly by combining precompositions into a factorization through $n-1$ variables and postcompositions into linear combinations. Lemma \ref{lem:lower} shows that $\Lambda$ vanishes on all such functions. Since $\Lambda(f)=1$, the clonoid $\mathcal C$ is not generated by its $(n-1)$-ary part. This completes the proof of Theorem \ref{thm:main}.

\section{Checks and further questions}

The proof above is independent of computation. The accompanying standard-library script \texttt{check.py} provides exact finite diagnostics of the Fourier projector definition, the nonzero-character sandwich, the zero-character quotient calculation, and the lower-bound functional. It uses arithmetic in finite prime fields, including targets containing the required roots of unity.

The lower-bound check includes all $262144$ matrices in $\Mat_3(\Z/4\Z)$ for $s=1$, and all $256$ matrices in $\Mat_2(\Z/4\Z)$ for $s=0$. It also checks one lift of every residue matrix in $\Mat_3(\F_3)$ and additional specified cases. These checks distinguish singular reduction from inner rank: the functional is tested on the larger set of matrices with singular reduction, which contains every matrix of the forbidden inner rank. They do not establish an infinite theorem by enumeration.

The method isolates why an elementary abelian factor costs one variable, while the cyclic $p^2$ factor can cost two: a nonzero congruence character is a rank-one matrix, and its nilpotent form requires two fixed coordinates. For sources with more than one cyclic $p^2$ factor, higher-rank congruence characters occur. The reduction in Lemma \ref{lem:rankone} then needs a different argument. Groups of higher exponent also require additional congruence layers. Neither case is settled here, and the general finite-module conjecture remains outside the scope of this paper.

\begin{thebibliography}{9}
\bibitem{FKR}
S. Fioravanti, M. Kompatscher and B. Rossi,
\emph{Clonoids over vector spaces}, preprint (2026),
\href{https://arxiv.org/abs/2602.04034v1}{arXiv:2602.04034v1}.

\bibitem{MW}
P. Mayr and P. Wynne,
\emph{Clonoids between modules},
International Journal of Algebra and Computation \textbf{34} (2024), 543--570.
\href{https://doi.org/10.1142/S021819672450022X}{doi:10.1142/S021819672450022X};
\href{https://arxiv.org/abs/2307.00076v3}{arXiv:2307.00076v3}.

\bibitem{Ko}
L. G. Kov\'acs,
\emph{Semigroup algebras of the full matrix semigroup over a finite field},
Proceedings of the American Mathematical Society \textbf{116} (1992), 911--919.
\href{https://doi.org/10.1090/S0002-9939-1992-1123658-6}{doi:10.1090/S0002-9939-1992-1123658-6}.

\bibitem{Kuhn}
N. J. Kuhn,
\emph{Generic representation theory of finite fields in nondescribing characteristic},
Advances in Mathematics \textbf{272} (2015), 598--610.
\href{https://doi.org/10.1016/j.aim.2014.12.012}{doi:10.1016/j.aim.2014.12.012};
\href{https://arxiv.org/abs/1405.0318v2}{arXiv:1405.0318v2}.
\end{thebibliography}
\end{document}
